\section*{第 \ref{ch:g12:destruction-of-mountains} 章——山脉的破坏}

\begin{solution}{exo:g12:destruction-of-mountains:1}
(a) \omterm{def:g12:destruction-of-mountains:denudation}{剥蚀速率}；(b) \omterm{def:g12:destruction-of-mountains:load}{推移质}；(c) \omterm{def:g12:destruction-of-mountains:load}{溶解质}；
(d) \omterm{def:g12:destruction-of-mountains:diagenesis}{成岩作用}；(e) \omterm{def:g12:destruction-of-mountains:rebound}{均衡回弹}。
\end{solution}

\begin{solution}{exo:g12:destruction-of-mountains:2}
(1) 砾岩。(2) 砂岩。(3) 页岩。(4) 角砾岩。
\end{solution}

\begin{solution}{exo:g12:destruction-of-mountains:3}
\qty{1}{m/Ma} 即 \num{1000000} 年内 \qty{1000}{mm}，
每年 \qty{0.001}{mm}。(1) 每年 \qty{0.054}{mm}；(2) 每年 \qty{0.218}{mm}。
(3) $\qty{10}{mm}$ 每年，即 1000 年内 \qty{10}{m}，即 \qty{10000}{m/Ma}。
\end{solution}

\begin{solution}{exo:g12:destruction-of-mountains:4}
固体负荷占总量的 \qty{85}{\percent}，因此总量为 $60/0.85
\approx 70.6$ 百万吨，\omterm{def:g12:destruction-of-mountains:load}{溶解质}负荷约为 $70.6 - 60 \approx
10.6$ 百万吨每年。
\end{solution}

\begin{solution}{exo:g12:destruction-of-mountains:5}
花岗岩的风化、河流的搬运、盆地中的沉积，
被较新地层掩埋、置换沉淀。
\end{solution}

\begin{solution}{exo:g12:destruction-of-mountains:6}
在 \qty{30}{cm/s} 时，沙粒被搬运（高于 \qty{20}{cm/s}）；
粉砂颗粒仍悬浮。在 \qty{5}{cm/s} 时，沙粒沉降
（低于 \qty{10}{cm/s}），但粉砂颗粒仍悬浮，因为
\qty{0.1}{cm/s} 即足以使其保持悬浮。
\end{solution}

\begin{solution}{exo:g12:destruction-of-mountains:7}
总负荷 $4/0.85 \approx 4.71$ 百万吨，$\num{4.71e9}/2830 \approx
\qty{1.66e6}{m^3}$ 的岩石（每年），分摊到 $\qty{5e10}{m^2}$ 上：$\num{3.3e-5}$ 米
每年，约 \qty{33}{m/Ma}，低于中位数 \qty{54}{m/Ma}。
\end{solution}

\begin{solution}{exo:g12:destruction-of-mountains:8}
(1) $1 - 2700/3300 \approx 0.182$。(2) 约 \qty{182}{m}。(3) 略多于
本章的 \qty{163}{m}：地壳相对地幔越轻，
柱体上升得越少。
\end{solution}

\begin{solution}{exo:g12:destruction-of-mountains:9}
(1) \qty{8}{km} 的岩石。(2) $8000 \times 0.163 \approx \qty{1300}{m}$。
(3) $8000/200 = 40$，即 4000 万年。
\end{solution}

\begin{solution}{exo:g12:destruction-of-mountains:10}
(1) $794/43 \approx 18$，即约 1800 万年。(2) 地表仅下降
被剥蚀厚度的 $0.163$ 倍，因此需移除的岩石厚度为 $794/0.163
\approx \qty{4900}{m}$，约需 $4900/43 \approx 110$，即约 1.1 亿年。
(3) 大陆不只是被磨蚀：造山作用与
岛弧的岩浆作用不断使其抬升并增生新地壳。
\end{solution}

\begin{solution}{exo:g12:destruction-of-mountains:11}
(1) $\num{1e14}/2830 \approx \qty{3.5e10}{m^3}$，即约
\qty{35}{km^3}。(2) $100/15.5 \approx 6.5$ 年。
\end{solution}

\begin{solution}{exo:g12:destruction-of-mountains:12}
由 $E = kh$ 与 $h = h_0 e^{-fkt}$ 可得 $E(t) = kh_0 e^{-fkt}$。
自 $t = 0$ 到 $t \to \infty$ 移除的厚度为 $\int_0^{\infty} kh_0
e^{-fkt}\,\mathrm{d}t = kh_0/(fk) = h_0/f$。要使地表降低 $h_0$，
剥蚀须移除 $h_0/f \approx 6.1\,h_0$，与\omterm{def:g12:destruction-of-mountains:rebound}{均衡回弹}命题所述一致：
即每剥蚀 1 米，地表仅降低 $f$ 米。
\end{solution}

\begin{solution}{exo:g12:destruction-of-mountains:13}
(1) $10 \times 1000 = \qty{10000}{mm}$，即 \qty{10}{m}。(2) 移去厚度为
$t$ 的冰层，相当于移去单位面积上质量为 $\rho_{\text{冰}}\,t$ 的
物质；柱体上升，直到它所排开的地幔减少了同样多：
$2000 \times 900/3380 \approx \qty{530}{m}$。(3) 地面随着接近新的平衡，
上升越来越少：驱动上升的是尚未
恢复的那一部分。
\end{solution}

\begin{solution}{exo:g12:destruction-of-mountains:14}
(1) 1900 年以前：约 \qty{56}{m/Ma}（如例中所示）。
1987--2006 年间：$145/0.85 \approx 171$ 百万吨，$\num{1.71e11}/2830
\approx \qty{6.0e7}{m^3}$，分摊到 $\qty{2.98e12}{m^2}$ 上，约
\qty{20}{m/Ma}。(2) 较早的那个更大：后期的负荷被水坝
等工程拦截。(3) 否：减少量为 $400 - 145 = 255$ 百万吨每
年，其中水坝占 100 至 150，约占一半。
\end{solution}

\begin{solution}{exo:g12:destruction-of-mountains:15}
(1) 至少 \qty{6}{km}。(2) $310 - 252 = 58$，即至多 5800 万年，因此
至少 $6000/58 \approx \qty{100}{m/Ma}$。(3) 其值介于中位数
（\qty{54}{m/Ma}）与各盆地平均值（\qty{218}{m/Ma}）之间，
这是造山带的常见速率。
\end{solution}

\begin{solution}{pb:g12:destruction-of-mountains:1}
\textbf{1.} $E = 0.209 \times 2000 \approx \qty{418}{m/Ma}$。

\textbf{2.} 每年 \qty{0.418}{mm}。

\textbf{3.} 比两者都快：这是一条陡峭的年轻山脉，速率与
地震活跃盆地相当，甚至更快。

\textbf{4.} 其平均坡度。

\textbf{5.} 厚 $\qty{0.418e-3}{m}$、覆盖 $\qty{1e6}{m^2}$ 面积的一层，
体积为 \qty{418}{m^3}，$418 \times 2830 \approx \qty{1.18e6}{kg}$：约 1200
吨。

\textbf{6.} $2000/418 \approx 4.8$ 百万年。

\textbf{7.} $h(t) = 2000\,e^{-0.209\,t}$（$t$ 以百万年为单位）。

\textbf{8.} $1/0.209 \approx 4.8$ 百万年。

\textbf{9.} $\ln 2/0.209 \approx 3.3$ 百万年。

\textbf{10.} $e^{-0.209 \times 11} = e^{-2.30} \approx 0.100$.

\textbf{11.} $1 - 2830/3380 \approx 0.163$.

\textbf{12.} 地表下降 $100 \times 0.163 \approx \qty{16}{m}$；
根部上升约 \qty{84}{m}。

\textbf{13.} $h' = -0.163 \times 0.209\,h \approx -0.0341\,h$，故 $h(t) =
2000\,e^{-0.0341\,t}$。

\textbf{14.} $(2000 - 200)/0.163 \approx \qty{11000}{m}$ 的岩石。

\textbf{15.} 曾位于约 \qty{11}{km} 深处的岩石：变质岩，
如片岩、片麻岩，以及花岗岩。

\textbf{16.} $1/0.0341 \approx 29$ 百万年。

\textbf{17.} $29 \times \ln 10 \approx 67$ 百万年。

\textbf{18.} $2000\,e^{-0.0341 \times 200} \approx 2000 \times 0.0011 \approx
\qty{2}{m}$。

\textbf{19.} $2000\,e^{-0.0341 \times 250} \approx 2000 \times 0.0002 \approx
\qty{0.4}{m}$：平坦的平原。山顶高度为 \qty{2037}{m}，说明
仅凭该模型不能拟合：该地区此后必定经历过再次隆升，或侵蚀速率
比模型假设的更慢。

\textbf{20.} 平均海拔减半所需时间为 $29 \times \ln 2 \approx$
\textbf{2000 万年}。
\end{solution}
